\chapter{反抗丘比}
每一个故事背后，至少都站着一位叙事者。
正如我在第五章所论，有理由把丘比（至少部分）视为《魔圆》的创作者，尤其是虚渊玄本人的化身。
因此，影片结尾虽有焰宣称自己已化为“邪恶”，恰如电视动画结尾的小圆化为“希望”，但片中更早出现的两场戏其实提示我们：
应当越过这一表象，找出真正的反派，也就是这场叛逆真正所要反抗的暴君。

影片开头不久，焰刚察觉魔法少女们的记忆遭人篡改，便发现囚禁她们的空间酷似魔女结界。
认出贝贝就是魔女夏洛特之后，她便认定这一切出自贝贝之手。
随后，小圆、麻美与焰围坐在麻美的餐桌旁。
那张桌子和电视动画里的一样，呈现为一个被拉长的等腰三角形。
小圆随口的一句评论，给了焰发问的机会：
她开始追问麻美，贝贝究竟从何而来，贝贝则显得困惑而不安。
整场演出里，焰都在验证一个猜想，即贝贝是这个密闭空间的创造者，亦即这个故事的反派。
然而，桌子的指向始终落在丘比身上：
无论是从上方俯瞰全体人物的高机位，还是更耐人寻味的沿桌面中线拍摄的镜头，都强调它越过贝贝，直指丘比。

与此相应，在影片接近结尾时，魔法少女们一开始看似是在与焰的使魔大军交战。
然而，当她们终于击碎结界，真正的景象才显露出来：
巨大的丘比从隔离场上方俯视着她们。
那才是真正的牢笼，困在其中的不只是其他魔法少女，也包括焰自己。
于是，我们又一次越过了眼前的焰，看见了藏在她背后的真正敌人。

上一章中，我曾简单谈到 E. T. A. 霍夫曼的《胡桃夹子与老鼠王》。
到了这里，这部作品值得进一步讨论，因为影片最后的大战，在许多地方都让人想起小说前半段玛丽的玩具军队与老鼠国王之间的那场战斗：
玩具们由一名断颌的怪物率领，从柜子里列队而出，与邪恶的鼠形怪物厮杀。
老鼠国王有七个头，头上戴着七顶王冠——
这些特征在原作中都有明确描写，因此很容易让人联想到《启示录》第十三章中那只有七个头、十顶冠的兽\footnote{%
    译注：七首十角的大红龙，就是魔鬼撒旦，意欲吞食身穿太阳的妇人所生之子，却在争战中被打落在地，后被天使捉住，以大链捆缚，丢入无底坑，禁锢一千年。
}。
而放到《魔圆》的语境里，“一个身体拥有多个头”这一形象，又恰好可以对应丘比的存在方式：
它拥有无数身体，却共享同一个意识。
正如老鼠国王与玛丽（Mary）为敌，《启示录》中的兽与上帝为敌一样，丘比也是小圆（Kannon）的敌人；
那么，焰作为一名为小圆而战的骑士，自然就对应着胡桃夹子。
最后，老鼠国王以饶过胡桃夹子为条件，要求玛丽不断交出自己的糖果和玩具供它吞食，这就是在索取赎金。
这和丘比要求魔法少女献出生命与灵魂、以满足它对能量的需求颇为相似。
毕竟，食物本来就是动物最基本的能量来源。
（这里还有一个值得注意的细节。
玛丽交出去的东西里，有一个可以吃的糖制人偶，那是她格外喜爱的玩具。
这个人偶的外貌，与恢复原貌后的胡桃夹子——
也就是现实中的德劳瑟迈耶那位侄子——
相当接近，尤其是两者都被形容为“面颊红润”。
如果胡桃夹子对应如今的焰，那么这尊面颊红润的糖人偶，便可以对应最初的焰：
那个体弱、无助、曾被小圆拯救、却也因此被她无意间卷入丘比轨道的少女。）

从这个角度来看，《胡桃夹子与老鼠王》中最耐人寻味的，或许正是德劳瑟迈耶这位叙事者身份上深深的不确定性。
德劳瑟迈耶把胡桃夹子送给玛丽，又亲口向她讲述胡桃夹子的身世以及他为何会受到诅咒。
然而，小说却屡屡暗示，他并不值得信任，甚至让人产生一种感觉：
他或许与老鼠国王存在某种隐秘的关联。
最明显的是，他任由时钟敲响，甚至也可以说正是他促成了钟声。
这招致老鼠国王入侵玛丽的房间，胡桃夹子也因此不得不召集玩具军队迎战。
后来，受伤的玛丽卧床休养时，德劳瑟迈耶又显得十分和蔼，给她讲故事，替她解闷。
可当玛丽告诉他，自己确实亲眼目睹玩具大战，后来也亲眼见证胡桃夹子杀死老鼠国王时，他却反过来把这一切斥为梦境和幻想。


更有意思的是，在讲述胡桃夹子诅咒来历的时候，德劳瑟迈耶还把“自己”写进了故事之中。
这与丘比进入焰的魔女结界有某种相似之处——
而它之所以能介入，恰恰又是因为结界从一开始就被包裹在丘比制造的隔离场之内。
按照德劳瑟迈耶自己的说法，在这则“故事中的故事”里，他差不多把自己塑造成了英雄：
为了替皮尔利帕特公主解除诅咒，他踏上旅程，四处寻找破解的方法。
然而，只要读者开始怀疑现实中的德劳瑟迈耶，内层故事中的这个德劳瑟迈耶就立刻可以被解读成一个经典童话中的邪恶叔父：
先是自己惹出麻烦，到了最后，却让自己的侄子替自己承担代价。
所谓解除诅咒的仪式真的只是“意外失败”了吗？
关于这一点，我们其实只有德劳瑟迈耶的一面之词。
甚至可以进一步怀疑：
也许那个仪式从一开始就不是为了彻底消除诅咒，而只是为了把诅咒从皮尔利帕特公主身上转移到胡桃夹子身上。
换句话说，德劳瑟迈耶讲述的故事就不可信。
与此同时，玛丽经历的那些奇异事件——
无论是玩具大战，还是后来进入胡桃夹子的奇幻国度——
都发生在一些很容易让人将其解释为梦境的时间和地点；
但在叙述方式上，与小说中那些毫无疑问属于日常“现实”的事件并没有本质区别。
而德劳瑟迈耶所讲的故事则不同，它从一开始就被明确框定为“讲给孩子听的、虚构的故事”。
因此，如果只从《胡桃夹子与老鼠王》自身的语境来判断，德劳瑟迈耶讲述的胡桃夹子身世，反倒比玛丽那些看似梦境的经历更缺乏“真实”性。
小说的结局甚至进一步暗示了这个论断：
因为最终，玛丽确实能够进入胡桃夹子的国度，而胡桃夹子本人说起德劳瑟迈耶时，态度也明显带着轻蔑。

当然，德劳瑟迈耶究竟是否真的比他表现出来的更加邪恶，最终仍取决于读者的解读。
不过，当玛丽终于进入胡桃夹子的国度，并使他恢复原貌时——
就像小圆把手伸进柜子，将焰从魔女状态中重新救回来一样——
胡桃夹子为什么会如此明确地贬低德劳瑟迈耶？
他告诉玛丽，德劳瑟迈耶根本无法实现自己曾经许诺的那片机械湖景；
相比之下，玛丽自己反而拥有比德劳瑟迈耶更强大的创造力量。
因此，故事最后给予玛丽的幸福结局是与胡桃夹子一起逃离这个“现实”，进入他真正所属的国度。

虽然具体情节，尤其事件顺序，无法严丝合缝地对应，人物之间的对应关系却非常明确。
焰创造了结界内的意象，反复与她相连的胡桃夹子主题说明，她把自己视作那个替公主承受诅咒、靠着终有一日能被玛丽救出这一希望而活下去的人物。
小圆既是公主，也是玛丽。
值得注意的是，在胡桃夹子的国度中，玛丽发现自己在水池中的倒影正是公主的模样；
在霍夫曼的故事中，两人很可能本就是同一角色。
丘比则兼具老鼠王与德劳瑟迈耶的特征。

当然，这些角色对应并不必然揭示人物的真实本性。
它们出自魔女结界的意象，而结界早已被确立为映照魔女心境、强加于现实并将其替换的幻觉迷宫。
但它们极大地揭示了焰的动机：
对玛丽而言，幸福的结局不是逃往某种涅槃或天堂，而是逃离现实、前往能与胡桃夹子共同生活的世界。
焰创造自己的世界，正是为了与属于她的玛丽，也就是小圆，在一起。
于是，她进行了一场邪恶而绝望的尝试，完成了属于她的胡桃夹子故事。
她比任何时候都更是那个阴暗丑陋的存在，而唯一的疗愈，是有人能够如其所是地爱她。
换言之，小圆必须了解焰做过什么、已经变成什么，却仍不放弃她。

但德劳瑟迈耶或许终会笑到最后。
作为故事里的叙事者，他必然在某种程度上映照着他自己的叙事者，即霍夫曼。
而在故事结尾，玛丽与胡桃夹子仍在霍夫曼的支配之中，因为只要他愿意，他就能继续书写他们，直至选择停笔。
丘比即德劳瑟迈耶，即霍夫曼，即虚渊玄。
它看似已经败北，在动画结局中也曾看似被打败，但最后那个凝视的眼睛提醒我们：
只要我们仍在观看，它也仍在观看。
只要我们愿意接受更多关于小圆、焰及其余所有人的故事，丘比就依旧拥有威胁她们的力量。
